\documentclass{article}
\usepackage{amsmath,amssymb,amsthm,hyperref}
\newtheorem{theorem}{Theorem}
\title{Synchronization by creation of lax limits in a native dialgebra doctrine}
\date{15 November 2026}
\begin{document}\maketitle
\paragraph{Status.} Written conditional general structural construction. This is an application of established enhanced/lax-algebra limit techniques, not a novelty assertion. It is not inferred from ordinary state pullbacks or from tight endpoint pullbacks of a double category. The strict 2-categorical hypotheses below are actual hypotheses, not an asserted strictification of the preceding arbitrary bicategorical setting.

\section{Independent input}
Let $L_0,K$ be strict 2-categories and $P,Q:L_0\to K$ strict 2-functors. Fix a small ordinary diagram shape $D$. Assume the selected $D$-shaped lax limits exist in $L_0$ and $Q$ preserves them as limits in $K$: applying $Q$ to a chosen universal lax cone gives a universal lax cone. No preservation assumption on $P$ is made. In formulas here $gf$ means first $f$, then $g$.

Define the 2-category $\operatorname{Dialg}_c(P,Q)$ as follows. Objects are $(X,f:PX\to QX)$. Arrows are
\[
(h,\alpha):(X,f)\to(Y,g),\qquad \alpha:Qh\,f\Rightarrow g\,Ph.
\]
Their composite uses $(kh,(\beta\,Ph)(Qk\,\alpha))$. A 2-cell $\theta:(h,\alpha)\Rightarrow(k,\beta)$ is $\theta:h\Rightarrow k$ such that
\[
(g\,P\theta)\alpha=\beta(Q\theta\,f).
\]
State maps and their 2-cells lie in $L_0$; transition cells lie in $K$. Composition and identities use their interchange laws. The subcategory of strict arrows has identity structural cells; these are the tight arrows. This orientation is specified explicitly rather than inferred from differing lax/colax naming conventions.

\section{Created synchronized transition}
Let a diagram in this doctrine have objects $(X_j,f_j)$ and arrows $(h_a,\alpha_a)$ for $a:j\to k$. Let the underlying chosen lax limit have projections $p_j:L\to X_j$ and coherent cells
\[
\tau_a:h_a p_j\Rightarrow p_k.
\]
The corresponding Cat-enriched weight is $W(j)=D/j$, with postcomposition on arrows; the identity arrows in these slices give the displayed generating projections. Here $\tau_1$ is identity and $\tau_{ba}=\tau_b(h_b\tau_a)$. The lax limit represents cones and their modifications by an isomorphism of hom categories. This strict representability is the universal property used, not just a bijection on objects.

Form a lax cone from $PL$ to the diagram $QX_j$ with legs $f_j Pp_j$ and comparison cells
\[
Qh_a\,f_jPp_j
\xRightarrow{\alpha_aPp_j}f_k Ph_a Pp_j
\xRightarrow{f_kP\tau_a}f_k Pp_k.
\]
The diagram-morphism coherence for $\alpha$ and lax-cone coherence for $\tau$, together with interchange, prove the cone equations. Since the $Qp_j,Q\tau_a$ cone is universal, it induces a unique
\[
f_L:PL\longrightarrow QL
\]
with $Qp_j f_L=f_jPp_j$ and $Q\tau_a f_L=(f_kP\tau_a)(\alpha_aPp_j)$. Thus each $p_j$ is a strict dialgebra arrow, and each $\tau_a$ is a dialgebra 2-cell. Agreement witnesses have not been discarded or inverted.

\begin{theorem}[Lax-limit creation, specified strict setting]
Under these assumptions, the forgetful 2-functor $\operatorname{Dialg}_c(P,Q)\to L_0$ creates the selected lax limits. Their projections are strict dialgebra arrows. An induced arrow is strict if its composites with all projections are strict.
\end{theorem}
\begin{proof}
Take any lax cone from $(Y,g)$ in the dialgebra doctrine. Its legs $(x_j,\gamma_j)$ have $\gamma_j:Qx_j g\Rightarrow f_j Px_j$. The underlying universal cone gives a unique $x:Y\to L$ with $p_jx=x_j$ and the prescribed cone cells. The family $\gamma_j$ is a modification between the two cones obtained from $Qx\,g$ and $f_L Px$; the required modification equation is exactly that the given cone comparison cells are dialgebra 2-cells. The hom-category universal property of the $Q$-image cone therefore gives a unique
\[
\gamma:Qx\,g\Rightarrow f_L Px
\]
whose projections are $\gamma_j$. This makes $(x,\gamma)$ the unique inducing dialgebra arrow. A modification of given cones induces a unique underlying $\theta:x\Rightarrow y$ by the lax limit of the $X_j$. Its dialgebra equation is true after all $Qp_j$, since it is true for each cone leg. Joint faithfulness on 2-cells of the universal $Q$-cone proves it in $K$. This establishes the required isomorphism of hom categories, including all noninvertible modifications. If all projected structural cells are identities, their modification equations first imply that the two projected cones have identical comparison cells as well as identical legs. Strict representability therefore gives equality $Qx\,g=f_LPx$ of the endpoints. Joint faithfulness can now compare $\gamma$ with the identity on that common endpoint, giving $\gamma=1$ and proving strictness detection. No native cell is assumed invertible.
\end{proof}

\section{Why this is not an ordinary state pullback}
For endotransitions $L_0=K$ and $P=Q=1_K$, the preservation condition is automatic. In $K=\mathbf{Cat}$, a cospan of systems $(X,f)\xrightarrow{(h,\alpha)}(Z,q)\xleftarrow{(k,\beta)}(Y,g)$ has lax-cone state objects
\[
(x,y,z,u:hx\to z,v:ky\to z).
\]
Its transition sends these to $(fx,gy,qz,q(u)\alpha_x,q(v)\beta_y)$, with the same formulas on morphisms. The retained arrows $u,v$ are synchronization data, not equality of the two projected states. This example merely checks the general formula; the theorem did not depend on category-valued states or on this cospan shape.

A strict pullback requires equality, while a pseudo-pullback retains an invertible comparison and cannot be substituted for this lax cone when $\alpha,\beta$ are directed. Conversely this lax synchronization does not assert that its comparison arrows represent physically legal sharing in every application. The admitted state/interface doctrine must decide that question.

\section{Relation to the existing native state presentation}
For fixed external $A,B$, the formal substitution $P(X)=X\otimes A$, $Q(X)=X\otimes B$ makes the displayed arrow cell exactly a native transition square, up to the earlier composition convention. This applies when these are actual strict 2-functors of a chosen state/process 2-category and the output action preserves the selected lax limits. A general relative presentation $R:S\to C\to H$ does not automatically satisfy that condition. In particular ordinary $S$ has no state 2-cells with which to retain a lax cone's coherence, and a nonfaithful relative pure map cannot be silently replaced by an inclusion.

Thus the theorem gives an independently motivated synchronization mechanism and a precise enrichment/preservation obligation. It does not establish the universal property inside the older ordinary native hom category simply by forgetting all modifications. The synchronized object and projection intertwiners may survive that forgetting, but the full lax-limit universal property belongs to the enriched doctrine stated above. Compatibility of these synchronized limits with the state-flattening pseudomonad, its pure interface and directed parallel assembly is a further theorem, not an automatic corollary.


\section{Changing-interface synchronization}
The two-base formulation matters. Suppose the state, input and output objects vary in a strict 2-category $K$ with selected strict 2-products and the selected lax limits. Let $R:K\to K$ preserve those limits and put
\[
L_0=K\times K\times K,\qquad
P(s,A,B)=Rs\times A,\qquad Q(s,A,B)=Rs\times B.
\]
A native arrow is now a triple $(h,a,b)$ with a transition square
\[
(Rh\times b)f\Rightarrow g(Rh\times a).
\]
Thus interface change is part of the same doctrine, rather than an afterthought restricted to fixed hom categories. The functor $Q$ preserves the selected lax limits: limits in $L_0$ are componentwise, $R$ preserves them, and the binary product 2-functor preserves weighted limits jointly. The last assertion follows directly from the hom-category product universal property, or from its being the right 2-adjoint of the diagonal. No assertion that product with one \emph{fixed} object preserves lax limits is involved.

The created synchronized process has the joint state, joint input and joint output given by their respective lax limits. Its transition and projection squares are obtained from the general theorem. In particular it retains output/interface agreement arrows instead of quietly treating them as equality. This is a general cartesian 2-categorical realization of the abstract two-base creation theorem, not a replacement of its hypotheses by a family of process examples.

For a noncartesian state/interface action the same construction applies if its output-action functor $Q$ preserves the selected limits jointly; cartesianity is sufficient here, not claimed necessary. For the original relative presentation this suggests an actual preservation/admissibility question: is its chosen state/interface doctrine closed under these enriched joint limits, and do the resulting maps remain admitted pure/native maps? Merely choosing a large ambient $K$ cannot answer that question for a restricted or nonfaithfully realized pure interface. No universal property is asserted after forgetting the native modifications.


\section{A relative tagged-interface application}
The heterogeneous form also prevents an unnecessary faithfulness assumption. Let $C$ be a strict cartesian 2-category, $R:C\to C$ and $J:C\to K$ strict 2-functors preserving the selected lax limits. No faithfulness of $J$ is required. Use
\[
L_0=C^3,\qquad P(s,A,B)=J(Rs\times A),\qquad Q(s,A,B)=J(Rs\times B).
\]
The same joint-product argument and the two preservation hypotheses prove that $Q$ preserves those lax limits. A native square retains its \emph{tagged} triple $(h,a,b)$ of actual maps in $C$, together with its transition cell in $K$. If $J$ identifies two tags, they are still distinct dialgebra arrows. Hence the creation theorem applies directly to this relative presentation without replacing its pure maps by their image or assuming an equipment with adjoints.

Take $I$ itself to be a chosen strict terminal object of $C$, and require $R(I)$ terminal. Thus state endomorphisms and state 2-cells at $I$ are unique. The unit-state sector has transitions $f:JA\to JB$, after transport through the actual product unitors and their $J$-images, and squares tagged by $a:A\to A'$, $b:B\to B'$ with $\alpha:Jb\,f\Rightarrow g\,Ja$. This recovers the relative-quintet square data, with the stated composition convention; enriched endpoint modifications are retained when $C$ has them. For locally discrete $C$, the corresponding ordinary tagged square category is recovered. Closure of this sector under synchronization additionally uses preservation of the terminal state and the chosen constant-state limit. Closure of the further pure-body marking $f=J(u)$ is a separate issue, not inferred from the square-data recovery.

This gives noncircular sufficient relative hypotheses: limits in the source doctrine, preservation by $R,J$, and the chosen product action. It is not a theorem that every original context/content factory has these properties, nor that an arbitrary bicategorical factory admits the assumed strict presentation. For other state actions the general criterion remains preservation by $Q$ itself, which may hold without Cartesianity or these factorwise sufficient conditions.

\section{Source-led placement}
Lack--Shulman, \emph{Enhanced 2-categories and limits for lax morphisms}, \url{https://arxiv.org/html/1104.2111v2}, explains why strict/tight projections and the variance of weak algebra morphisms are essential to limit lifting. The direct proof above spells out this mechanism for a native dialgebra presentation; no application to an arbitrary pseudomonad is inferred without a bridge. There is an exact established bridge for endotransitions in $\mathbf{Cat}$: strict algebras of the 2-monad $T(X)=\mathbb N_{\mathrm{disc}}\times X$ are categories with a chosen endofunctor, and a colax algebra map is determined by its degree-one cell $hf\Rightarrow gh$. Higher components are its iterated pastes. Section 3.5.2 of Lack--Shulman explicitly treats the corresponding colax-algebra lax-limit orientation, with tight generating projections and joint tightness detection. Thus that endotransition case is an instance of existing general machinery, not a newly discovered feedback structure.

Loregian, \emph{Monads and limits in bicategories of circuits}, \url{https://arxiv.org/html/2501.01882}, Theorem 30 gives tight endpoint pullbacks in the stated Mealy setting, while Section 4 uses inserter/comma/pullback descriptions. These are relevant established machinery, but Theorem 30 alone is not the machine-hom synchronization assertion proved here. A source-identification check for the exact general dialgebra formulation remains open; no novelty claim is made.
\end{document}
